\scanpage{417}
\mspage{330}
% archive sheet number: 416

% 1–3
1) $\varphi$ et $\psi$ sont cofinaux, $\psi$ est l'inclusion d'une sous-catégorie str. pleine $J'$ de $J$, $I$ est filtrant (donc $J$ et $J'$ aussi).

% 4–10
On cherche des conditions pour que l'inclusion \deleted{\unsure{pl.}} de la sous-catégorie str. pleine $I'\unsure{-|\to}I$ soit cofinale (auquel cas on voit tout de suite que $I'$ est \unsure{filtrante|filtrant}, et $I'\to J'$ est aussi cofinal). Pour ceci, on suppose qu'il existe un \uline{ordinal limite} $\alpha_0$ ayant les propriétés \inserted{suppl.} suivantes\unsure{.|:}

% 11–13
2) $I$ et $J'$ sont stables par $\varinjlim$ du type $I_{\alpha_0}$ $(=\{\beta\mid\beta\ \text{ordinal}<\alpha_0\})$, et les foncteurs $\varphi,\psi$ commutent à ces limites.

% 14–18
3) Pour tt $\alpha\in I_{\alpha_0}$, $I$ est \deleted{\unsure{stab}} \unsure{grd|grand} devant $\operatorname{Card}\alpha$. [Donc, $\varphi$ étant cofinal, $J$ est \inserted{\unsure{égalemt|églmt|également}} grd devant $\operatorname{Card}\alpha$, et de $=$, $\psi$ étant pl. fidèle, et cofinal, $J'$ est grd devant $\operatorname{Card}\alpha$.]

% 19–20
Sous ces conditions, $I'$ est cofinal dans $I$ (donc \inserted{aussi} grand devant $\operatorname{Card}\alpha$, pour tt $\alpha\in I_{\alpha_0}$).

% 21–25
La démonstration est celle donnée \deleted{dans} ci-dessus. On n'a \unsure{pas eu} à utiliser que l'\deleted{\unsure{entier}}ordinal $\alpha_0=\alpha_{\unsure{\pi|\varpi}}$ \unsure{était} de cette espèce particulière, \unsure{seulement} que c'était un ordinal limite.
